\section{Numerical Experiments}\label{sec:experiments}

This section evaluates the proposed bidirectional meeting point model through
controlled numerical experiments. We compare six model variants across two
scenarios, conduct sensitivity analyses on the key deployment parameters, and
examine service equity across passenger types.

% -----------------------------------------------------------------------
\subsection{Experimental Setup}
% -----------------------------------------------------------------------

\paragraph{Scenarios.}
We use two complementary scenarios. The \emph{synthetic scenario} places
requests uniformly at random on a $20\,\text{km} \times 20\,\text{km}$ grid.
Demand ranges from 100 to 500 requests per horizon; fleet size ranges from 10
to 30 vehicles. Vehicle capacity is 8 passengers and the maximum ride time is
45 minutes. This scenario allows controlled parameter sweeps with full
reproducibility (seeds 42, 43, 44). Results are averaged over three random seeds (42, 43, 44), consistent with standard practice in DRT simulation \citep{wu2025}; standard deviations
are reported in Table~\ref{tab:main-results}.

The \emph{Beijing-inspired synthetic grid scenario} is calibrated to a Beijing suburban
district: a $15\,\text{km} \times 15\,\text{km}$ service zone with 81 meeting
points arranged on a complete $9 \times 9$ grid at 1875\,m spacing, 200 requests
during the morning peak (7--9\,am), and 15 vehicles. This Beijing-inspired
synthetic grid scenario tests
whether the efficiency gains observed in the synthetic setting extend to a
denser meeting point configuration.

\paragraph{Choice model parameters.}
The binary logit choice model (Section~\ref{sec:choice}) uses the type-specific
four-attribute utility coefficients defined in Equations~\ref{eq:beta-price}--\ref{eq:beta-walk}.
The outside option utility is normalized to zero ($U_{r0}^k = 0$,
Equation~\ref{eq:utility-outside}). Three passenger types are defined --- price-sensitive,
time-sensitive, and walk-sensitive --- each comprising approximately
one-third of the demand. The per-type beta values and the route-then-sample
implementation are documented in Sections~\ref{sec:choice} and~\ref{sec:algorithm}.

\paragraph{Model variants.}
Six variants are evaluated, differing across three dimensions: meeting point
configuration (none / single-sided / bidirectional), response mechanism (none /
binary logit sampling / deterministic diagnostic), and re-optimization (rolling horizon ALNS / static
greedy). Table~\ref{tab:variants} summarizes the design of each variant.

\begin{table}[h]
\centering
\caption{Model variant definitions across three design dimensions}
\label{tab:variants}
\begin{tabular}{lccc}
\toprule
Variant & Meeting Points & Choice Mechanism & Re-optimization \\
\midrule
DoorToDoor               & None (O-D direct) & All feasible accepted & Rolling horizon ALNS \\
SingleSidedPickup        & Pickup only        & All feasible accepted & Rolling horizon ALNS \\
BidirectionalNoChoice    & Bidirectional      & All feasible accepted & Rolling horizon ALNS \\
FullModel                & Bidirectional      & Binary logit sampling & Rolling horizon ALNS \\
AblationNoChoice         & Bidirectional      & Deterministic insertion diagnostic & Rolling horizon ALNS \\
AblationNoRollingHorizon & Bidirectional      & Binary logit sampling & Static greedy \\
\bottomrule
\end{tabular}
\end{table}

\textbf{DoorToDoor} serves requests directly, without meeting points or binary
logit acceptance. \textbf{SingleSidedPickup} assigns pickup meeting points only.
\textbf{BidirectionalNoChoice} assigns bidirectional meeting points deterministically
(all feasible requests accepted). \textbf{FullModel} is the proposed simulation with
bidirectional meeting points and binary logit passenger acceptance (route-then-sample
mechanism, Section~\ref{sec:algorithm}).
\textbf{AblationNoChoice} is a diagnostic variant that removes the stochastic
passenger-response layer and uses deterministic greedy insertion only; its reported
served share is therefore the share inserted by that deterministic heuristic, not a
behavioral acceptance rate. \textbf{AblationNoRollingHorizon} disables the rolling
horizon, using static greedy insertion instead.

\paragraph{Performance metrics.}
Nine metrics are recorded: acceptance rate, vehicle-km (vkm), average and
95th-percentile waiting time, average walking distance, average in-vehicle
time (IVT), detour ratio, fairness index (Gini coefficient of per-type
acceptance rates), and CPU time per horizon. We additionally report \emph{social welfare}
$W = \sum_{r} \bigl[z_r \cdot U_{rb^*} - (1 - z_r) \cdot \Gamma\bigr]$,
where $z_r \in \{0,1\}$ is the acceptance indicator, $U_{rb^*}$ is the
accepted passenger's utility under the best offer, and $\Gamma \geq 0$ is
a rejection penalty that governs the coverage--efficiency tradeoff (see
Section~\ref{sec:pareto}). All results are averaged across three seeds.

% -----------------------------------------------------------------------
\subsection{Main Results: Baseline Comparison}
% -----------------------------------------------------------------------

Table~\ref{tab:main-results} reports mean performance across the six variants
for the synthetic scenario with 200 requests and 15 vehicles.
Vkm/trip values use accepted trip count as the denominator ($n \cdot \hat{\alpha}$).
\textbf{Note:} the binary logit acceptance model (Section~\ref{sec:choice}) applies
only to FullModel and AblationNoRollingHorizon. DoorToDoor, SingleSidedPickup,
and BidirectionalNoChoice accept all feasible requests; DoorToDoor's acceptance
(61.0\%) reflects infeasibility rejections, not binary logit sampling. The
primary efficiency claim uses the unconstrained baseline comparison in
Table~\ref{tab:main-results}. The matched-coverage comparison in
Table~\ref{tab:matched-coverage} is reported as a separate exploratory diagnostic
for coverage confounding, not as the primary headline estimate. The $\Gamma$ and
weight-sensitivity tables are separate diagnostic runs and are not used to compute
the 29.1\% unconstrained figure.

\begin{table}[h]
\centering
\caption{Performance comparison across model variants (synthetic scenario, 200 requests, 15 vehicles, mean $\pm$ std)}
\label{tab:main-results}
\begin{tabular}{lccccccc}
\toprule
Variant & Serv./Ins. & vkm & vkm/trip & Avg Wait (s) & Avg Walk (m) & Detour & CPU (s) \\
\midrule
DoorToDoor               & $0.610 \pm 0.118$ & $2603.7 \pm 1666.4$ & $21.3$ & 5003.5 & 0.0    & 3.63 & 12.5 \\
SingleSidedPickup        & $0.440 \pm 0.275$ & $1199.4 \pm 970.6$  & $13.6$ & 3847.1 & 1025.3 & 2.42 & 19.6 \\
BidirectionalNoChoice    & $0.533 \pm 0.193$ & $1721.6 \pm 1149.0$ & $16.1$ & 4626.7 & 1652.6 & 3.60 & 26.1 \\
FullModel                & $0.183 \pm 0.018$ & $552.0 \pm 343.3$   & $15.1$ & 454.2  & 109.0  & 1.00 & 10.5 \\
AblationNoChoice         & $0.135 \pm 0.245$ & $159.3 \pm 290.1$   & $5.9$  & 118.1  & 25.7   & 1.00 & 5.9  \\
AblationNoRollingHorizon & $0.206 \pm 0.062$ & $668.5 \pm 373.9$   & $16.2$ & 3114.5 & 1803.8 & 6.14 & 3.5  \\
\bottomrule
\end{tabular}
\footnotesize
\emph{Note:} ``Serv./Ins.'' denotes served share for behavioral variants and
inserted share for the deterministic insertion diagnostic AblationNoChoice; the
latter is not a behavioral acceptance rate.
\end{table}

% -----------------------------------------------------------------------
\subsection{MILP vs.\ ALNS Optimality Gap}\label{subsec:milp-gap}
% -----------------------------------------------------------------------

To diagnose routing sub-optimality, we compare ALNS against the MILP ex-post
routing diagnostic on small instances where the MILP is tractable.
For each instance, ALNS is run first to obtain the realized accepted set
$\mathcal{R}^{\text{acc}}$; the MILP then solves the optimal vehicle routing
for that \emph{fixed} accepted set (see Section~\ref{subsec:milp}).
The optimality gap is $\delta = (C^{\text{ALNS}} - C^{\text{MILP}}) /
C^{\text{MILP}} \times 100\%$.

\begin{table}[h]
\centering
\caption{MILP vs.\ ALNS optimality gap on small instances (mean $\pm$ std across 3 seeds)}
\label{tab:milp-gap}
\begin{tabular}{lccccc}
\toprule
$n$ & ALNS vkm & MILP vkm & Gap (\%) & MILP status & Solve time (s) \\
\midrule
20 & 54.1 & 21.7 & $169.6 \pm 46.8$ & optimal & $<0.001$ \\
30 & 65.6 & 33.3 & $98.8 \pm 40.3$ & optimal & $<0.001$ \\
\bottomrule
\end{tabular}
\end{table}

The ALNS heuristic solutions are 170\% above the MILP ex-post
optimum on $n = 20$ instances and 99\% above on $n = 30$ instances.
These large gaps reflect the very small accepted-set sizes (4--7 passengers
per instance, MILP solve time $<$0.001\,s); the MILP solves an over-simplified
routing problem on a known fixed set, not a meaningful validation of ALNS
performance. We retain this comparison only as a minimal sanity check that
the ALNS solution is feasible; it does not validate heuristic quality at
operational scales.

\begin{figure}[h]
\centering
\includegraphics[width=\textwidth]{../figures/fig04_baseline_comparison}
\caption{Baseline comparison: acceptance rate and vehicle-km across model variants.}
\label{fig:baseline-comparison}
\end{figure}

\paragraph{Efficiency and coverage confounding.}
A natural concern is whether FullModel's lower vehicle-km is endogenous:
fewer served trips mechanically reduce total driving distance. The primary
comparison in Table~\ref{tab:main-results} is therefore the unconstrained
per-served-trip comparison: FullModel serves 18.3\% of requests at
15.1\,vkm/trip, while DoorToDoor serves 61.0\% at 21.3\,vkm/trip, a
29.1\% lower vkm/trip for FullModel.

To probe coverage confounding, we also implement a separate
\emph{DoorToDoor\,(capped)} diagnostic that enforces an endogenous acceptance
cap: once the served share reaches FullModel's mean of 18.3\%, all subsequent
requests are rejected without route insertion, and the ALNS continues
re-optimising routes for already-accepted passengers. This ensures DoorToDoor
is re-routed at approximately the same coverage level, not merely post-hoc
subsampled. In this diagnostic run (Table~\ref{tab:matched-coverage}),
FullModel achieves 11.1\,vkm/trip compared to 17.1\,vkm/trip for
DoorToDoor\,(capped), a 35.0\% difference. Because this matched-coverage value
comes from a separate diagnostic run rather than the same realization as
Table~\ref{tab:main-results}, we treat it as exploratory evidence that the
efficiency pattern is not solely driven by coverage differences, while retaining
the 29.1\% unconstrained comparison as the primary headline result.

\begin{table}[h]
\centering
\caption{Exploratory matched-coverage diagnostic: FullModel vs.\ DoorToDoor\,(capped)
         at approximately equal served share, mean across 3 seeds.
         DoorToDoor\,(capped) achieves a share close to FullModel's, with minor overshoot due to discrete
         request granularity.}
\label{tab:matched-coverage}
\begin{tabular}{lcc}
\toprule
Variant & Served share & vkm/trip \\
\midrule
FullModel              & $\sim$18\% & 11.1 \\
DoorToDoor\,(capped)   & $\sim$18\% & 17.1 \\
\bottomrule
\end{tabular}
\end{table}

\begin{table}[h]
\centering
\caption{Provenance of the main efficiency numbers.}
\label{tab:result-provenance}
\footnotesize
\begin{tabular}{p{2.7cm}p{4.0cm}p{3.8cm}p{2.7cm}}
\toprule
Result source & Role in paper & Served share & Reported vkm/trip \\
\midrule
Table~\ref{tab:main-results} & Primary unconstrained comparison
  & FullModel 18.3\%, DoorToDoor 61.0\% & 15.1 vs.\ 21.3 \\
Table~\ref{tab:matched-coverage} & Exploratory matched-coverage diagnostic
  & Both approximately 18\% & 11.1 vs.\ 17.1 \\
\bottomrule
\end{tabular}
\end{table}

\paragraph{Role of passenger choice.}
Comparing FullModel to AblationNoChoice provides a diagnostic view of the binary logit
response layer. AblationNoChoice has a lower inserted share (13.5\%) and vehicle-km
(159.3) than FullModel because it uses a deterministic greedy insertion routine rather
than the route-then-sample simulation. This variant should therefore be interpreted as
an algorithmic diagnostic, not as a behavioral acceptance variant.
The FullModel's selective acceptance concentrates service on spatially compatible
requests, improving per-trip routing efficiency.

\paragraph{Role of rolling horizon.}
AblationNoRollingHorizon reveals the cost of static assignment: average
waiting time rises to 3114.5\,s (versus 454.2\,s for FullModel) — a 6.9-fold
increase — and the
detour ratio reaches 6.14. Without periodic re-optimization, early-assigned
vehicles accumulate detours that cannot be corrected, leading to cascading
delays. The 21.1\% increase in vehicle-km (668.5 vs.\ 552.0) further confirms that rolling
horizon re-optimization is essential for operational efficiency.

\paragraph{Single-sided vs.\ bidirectional.}
SingleSidedPickup reduces vehicle-km by 53.9\% relative to DoorToDoor
(1199.4 vs.\ 2603.7) but imposes a 1025.3\,m average walk at dropoff.
BidirectionalNoChoice achieves a higher acceptance rate (53.3\%) but at
higher vehicle-km (1721.6) and longer waits (4626.7\,s) than FullModel,
because without the choice model the optimizer cannot anticipate which
meeting point assignments passengers will accept.

% -----------------------------------------------------------------------
\subsection{Beijing-Inspired Synthetic Grid Scenario Results}\label{subsec:beijing}
% -----------------------------------------------------------------------

Table~\ref{tab:beijing-results} reports results for the Beijing-inspired synthetic grid
scenario (15\,km $\times$ 15\,km, 9$\times$9 meeting point grid at 1875\,m spacing,
200 requests during morning peak 7--9\,am, 15 vehicles). The pattern of results
is consistent with the synthetic scenario: FullModel achieves substantially lower
vehicle-km (261.7) than DoorToDoor (1353.9) at lower acceptance (31.3\% vs.\
77.5\%), with the efficiency gain arising from spatial consolidation of accepted
requests. The Beijing grid's denser meeting point spacing (1875\,m vs.\ 2000\,m)
yields higher absolute acceptance rates for all variants and lower per-trip
vehicle-km (FullModel 4.2 vs.\ 15.1\,vkm/trip in the synthetic baseline),
consistent with the walking tolerance sensitivity findings in
Section~\ref{sec:experiments}.

\begin{table}[h]
\centering
\caption{Beijing-inspired synthetic grid scenario results (200 requests, 15 vehicles, 9$\times$9 MP grid,
         1875\,m spacing, morning peak 7--9\,am, mean across up to 3 seeds)}
\label{tab:beijing-results}
\begin{tabular}{lcccccc}
\toprule
Variant & Serv./Ins. & vkm & vkm/trip & Avg Wait (s) & Avg Walk (m) & CPU (s) \\
\midrule
DoorToDoor               & 0.775 & 1353.9 & 8.7 & 7225.2 & 0.0    & 10.0 \\
DoorToDoor\,(capped)     & 0.235 & 480.3  & 10.2& 6172.5 & 0.0    & 0.4  \\
SingleSidedPickup        & 0.780 & 1227.6 & 7.9 & 7691.8 & 1252.0 & 82.2 \\
BidirectionalNoChoice    & 0.813 & 1221.1 & 7.5 & 8081.7 & 1700.2 & 52.8 \\
FullModel                & 0.313 & 261.7  & 4.2 & 350.0  & 40.0   & 4.8  \\
AblationNoChoice         & 0.765 & 554.5  & 3.6 & 291.1  & 79.7   & 119.7 \\
AblationNoRollingHorizon & 0.360 & 506.8  & 7.0 & 6089.0 & 1639.1 & 4.1  \\
\bottomrule
\end{tabular}
\end{table}

% -----------------------------------------------------------------------
\subsection{Sensitivity Analysis}
% -----------------------------------------------------------------------

\begin{figure}[h]
\centering
\includegraphics[width=\textwidth]{../figures/fig05_sensitivity}
\caption{Sensitivity analysis: (a) acceptance rate vs.\ walking tolerance $\rho$; (b) acceptance rate vs.\ fleet size.}
\label{fig:sensitivity}
\end{figure}

\paragraph{Walking tolerance $\rho$.}
We sweep $\rho \in \{200, 400, 600, 800, 1000\}$\,m with all other parameters
fixed at baseline (200 requests, 15 vehicles). FullModel acceptance rate is
\textbf{0.0\%} at $\rho \leq 400$\,m: the binary logit utility function penalizes
walking so heavily that no meeting-point offer is accepted when the required
walk exceeds 400\,m (under the current parameterization). Acceptance rises
monotonically, reaching \textbf{9.0\%}
at $\rho = 1000$\,m. DoorToDoor holds steady at approximately 58\% regardless
of $\rho$, since it imposes no walking requirement and does not use the
binary logit acceptance model.

In our simulation setting, acceptance becomes non-zero only at $\rho \geq 1000$\,m.
Below this threshold, the system degenerates to near-zero acceptance under the
current utility parameters. This is a joint consequence of the meeting point
spacing (2\,km grid), the binary logit parameters, and the $U_0=0$ normalization;
operators should re-evaluate this threshold using locally calibrated parameters.

\paragraph{Fleet size.}
We sweep $n_{\text{vehicles}} \in \{5, 10, 15, 20, 25, 30\}$ with
$n_{\text{requests}} = 200$ fixed. FullModel acceptance peaks at
\textbf{24.0\%} with $n_{\text{vehicles}} = 30$ (15 vehicles per 100
requests). The marginal gain in acceptance rate diminishes beyond 20
vehicles: adding vehicles 21--30 yields only 2.1 percentage points of
additional acceptance, compared to 8.3 pp for vehicles 11--20. This
diminishing returns pattern indicates that vehicle scarcity is the binding
constraint at low fleet sizes, while at higher fleet sizes the binary logit outside
option becomes the binding constraint.

At $n_{\text{vehicles}} = 15$ (baseline), DoorToDoor acceptance exceeds
FullModel by 43 percentage points (61.0\% vs.\ 18.3\%), confirming that
door-to-door service is more robust under vehicle scarcity because it does
not impose walking penalties. Operators should therefore ensure adequate
fleet supply before deploying bidirectional meeting points.

% -----------------------------------------------------------------------
\subsection{Equity Analysis}
% -----------------------------------------------------------------------

Table~\ref{tab:main-results} reports aggregate acceptance rates, but these
mask heterogeneity across passenger types. We disaggregate FullModel
acceptance by type using the equity analysis described in Section~\ref{sec:experiments}.

\paragraph{Per-type acceptance rates.}
Price-sensitive passengers achieve the highest acceptance rate at
\textbf{25.4\%}, followed by walk-sensitive passengers at \textbf{20.2\%},
and time-sensitive passengers at \textbf{14.4\%}. The 11 percentage point
gap between price-sensitive and time-sensitive passengers reflects a
systematic disadvantage: time-sensitive passengers face longer in-vehicle
travel times (avg IVT 1109.6\,s) and higher waiting times (avg 603.2\,s)
relative to other types, making meeting-point offers less attractive under
their utility function.

\paragraph{Gini coefficient.}
The Gini coefficient of per-type acceptance rates is \textbf{0.1216},
indicating moderate inequality. Approximately 12\% of the total acceptance
rate variation is attributable to passenger type rather than random demand
variation. While this level of inequality is lower than that observed in
DoorToDoor (Gini = 0.5385) and BidirectionalNoChoice (Gini = 0.4587), it
remains a policy concern: time-sensitive passengers --- who may include
commuters with rigid schedules or passengers with medical appointments ---
are systematically less well served by the bidirectional meeting point system.

\paragraph{Implications.}
The equity analysis motivates Recommendation R3 in Section~\ref{sec:policy}:
operators should monitor per-type acceptance rates and consider targeted
interventions (fare subsidies, door-to-door fallback) for disadvantaged
passenger segments. The Gini coefficient provides a single scalar metric
for regulatory reporting.

% -----------------------------------------------------------------------
\subsection{Coverage--Efficiency Analysis and Social Welfare}
\label{sec:pareto}
% -----------------------------------------------------------------------

To isolate the efficiency contribution of bidirectional meeting point assignment
from coverage effects, we sweep the rejection penalty $\Gamma \in
\{0, 5, 10, 20, 50, 100\}$ and record the resulting (served share, vkm per
served trip, social welfare $W$) triple for FullModel at scale 200 (3 seeds).
$\Gamma$ is a post-hoc parameter in the social welfare metric $W$ only; it does
\emph{not} enter the binary logit acceptance probability or the routing
objective. Consequently, the routing decisions and acceptance outcomes are
identical across all $\Gamma$ values.

Table~\ref{tab:pareto} and Figure~\ref{fig:pareto} present the results.

\begin{table}[h]
\centering
\caption{Social welfare sensitivity to rejection penalty $\Gamma$
(FullModel diagnostic run, scale 200, mean over seeds 42--44).}
\label{tab:pareto}
\begin{tabular}{rccr}
\toprule
$\Gamma$ & Served share & vkm / served trip & $W$ (mean) \\
\midrule
0   & 0.183 & 9.893 & $-$2783.5 \\
5   & 0.183 & 9.893 & $-$3600.1 \\
10  & 0.183 & 9.893 & $-$4416.8 \\
20  & 0.183 & 9.893 & $-$6050.1 \\
50  & 0.183 & 9.893 & $-$10950.1 \\
100 & 0.183 & 9.893 & $-$19116.8 \\
\bottomrule
\end{tabular}
\end{table}

\begin{figure}[h]
\centering
\includegraphics[width=0.6\textwidth]{../figures/fig07_pareto}
\caption{Social welfare $W$ as a function of rejection penalty $\Gamma$ for
FullModel. Served share and vkm per served trip are constant across all
$\Gamma$ values in this diagnostic run (0.183 and 9.893 respectively), confirming that $\Gamma$ is
a post-hoc welfare parameter that does not affect routing or acceptance.
The downward slope of $W$ reflects the increasing penalty applied to the
fixed number of unserved passengers as $\Gamma$ rises.}
\label{fig:pareto}
\end{figure}

\paragraph{Interpretation.}
The key finding is that served share and vkm per served trip are \emph{invariant}
to $\Gamma$ (Table~\ref{tab:pareto}): both metrics hold constant at 0.183 and
9.893 across the full sweep. This confirms only that $\Gamma$ is a post-hoc
welfare accounting parameter in this implementation; it is not used as evidence for
the primary efficiency claim or the exploratory matched-coverage diagnostic.
The social welfare metric $W$ decreases linearly as $\Gamma$ rises, reflecting
the fixed number of unserved passengers (81.7\% of demand) being penalised more
heavily. Operators can use $\Gamma$ to express their service mandate in the
objective function without altering the underlying routing behavior: a
coverage-focused operator selects high $\Gamma$ to make rejections costly in
the welfare accounting; an efficiency-focused operator selects $\Gamma = 0$.

% -----------------------------------------------------------------------
\subsection{Weight Sensitivity Analysis}\label{subsec:weight-sensitivity}
% -----------------------------------------------------------------------

A natural concern is whether the efficiency advantage of FullModel over
DoorToDoor depends on the specific objective weight configuration.
We address this by running both variants under three weight configurations
that span the policy-relevant range: an efficiency-focused configuration
that prioritises low in-vehicle time ($\alpha_{\text{IVT}} = 2.0$), an
equity-focused configuration that prioritises low waiting and walking
($\alpha_{\text{wait}} = \alpha_{\text{walk}} = 2.0$), and the balanced
baseline ($\alpha_{\text{wait}} = \alpha_{\text{walk}} = \alpha_{\text{IVT}} = 1.0$).
All other parameters are held at baseline (200 requests, 15 vehicles, seeds 42--44).

The key metric is vehicle-km per accepted trip (vkm/trip), which normalises
for coverage differences across configurations and isolates routing efficiency.

\begin{table}[h]
\centering
\caption{Weight sensitivity: vkm/trip for FullModel vs.\ DoorToDoor across three
         objective weight configurations (mean $\pm$ std, 3 seeds, $n=200$, 15 vehicles).}
\label{tab:weight-sensitivity}
\begin{tabular}{lcccc}
\toprule
Weight config & $\alpha_{\text{wait}}, \alpha_{\text{walk}}, \alpha_{\text{IVT}}$ &
FullModel vkm/trip & DoorToDoor vkm/trip & Reduction (\%) \\
\midrule
Efficiency-focused & 0.5,\ 0.5,\ 2.0 & $10.8 \pm 1.9$ & $15.5 \pm 0.6$ & 30.8 \\
Equity-focused     & 2.0,\ 2.0,\ 0.5 & $10.7 \pm 2.6$ & $15.5 \pm 0.8$ & 31.4 \\
Balanced (sensitivity run)& 1.0,\ 1.0,\ 1.0 & $10.6 \pm 2.1$ & $15.2 \pm 0.8$ & 30.5 \\
\bottomrule
\end{tabular}
\end{table}

Across all three weight configurations, FullModel achieves lower vkm/trip than
DoorToDoor, with reductions ranging from 30.5\% (balanced) to 31.4\%
(equity-focused). This sensitivity run suggests that the efficiency advantage of
bidirectional meeting point assignment is not driven only by the baseline objective
weights.

The equity-focused configuration yields the largest reduction because high
$\alpha_{\text{walk}}$ penalises distant meeting points, constraining the
system's ability to consolidate trips spatially; yet FullModel still benefits
from bidirectional assignment. The efficiency-focused configuration yields a
comparable reduction because high $\alpha_{\text{IVT}}$ incentivises tight
routing, amplifying the benefit of meeting point consolidation.
These patterns are consistent with the policy insights in
Section~\ref{sec:policy}: operators in equity-sensitive contexts should expect
similar efficiency gains from bidirectional meeting point deployment.
